% <!-- PRIMARY-SOURCE-EXEMPT: reason=paper with existing bibliography (chbmit/Goldberger 2000) -->
\documentclass[11pt]{article}

\usepackage{amsmath,amssymb,amsthm}
\usepackage{booktabs}
\usepackage[margin=2.5cm]{geometry}
\usepackage{hyperref}
\usepackage{enumitem}

\newtheorem{remark}{Remark}[section]

\begin{document}

\title{Discrete Fibonacci Orbit Classification Adds $\Delta R^2 = +0.21$\\
to EEG Seizure Prediction Beyond Spectral Baselines:\\
A 10-Patient CHB-MIT Study}
\author{Will Dale\\
\small Independent Researcher}
\date{}
\maketitle

\begin{abstract}
We report that features derived from the Fibonacci orbit decomposition of
modular arithmetic --- applied as a discrete observer projection to topographic
EEG clustering --- add mean $\Delta R^2 = +0.210$ to seizure-period prediction
beyond a delta-band spectral baseline in 10 patients from the CHB-MIT Scalp EEG
Database (Fisher combined $\chi^2 = 208$, $p = 2.9 \times 10^{-33}$;
9/10 patients individually significant).
The algebraic structure (Quantum Arithmetic at modulus $m=9$) is applied
without any parameter fitting: a fixed discretisation map converts spectral
centroids to modular state pairs $(b,e)$, producing 4-tuple features
$(b,e,d,a)$ where $d=(b+e-1)\bmod 9 + 1$ and $a=(b+2e-1)\bmod 9 + 1$.
These four integers augment a standard regression model.
We describe the experimental design, the observer-projection firewall
that prevents circular reasoning, the per-patient results, and the
limitations that a pre-registered replication study should address.
\end{abstract}

% ====================================================================
\section{Introduction}\label{sec:intro}

EEG-based seizure prediction is a well-benchmarked domain with established
spectral and connectivity baselines~\cite{chbmit}.  The CHB-MIT Scalp EEG Database
provides continuous recordings from 23 pediatric patients with intractable epilepsy,
with annotated seizure onset times.  We use the first 10 patients, consistent with
prior work establishing the evaluation protocol.

The Quantum Arithmetic (QA) system operates on pairs $(b,e) \in \{1,\dots,m\}^2$
under the Fibonacci shift $T(b,e) = (e,\, ((b+e-1)\bmod m)+1)$.
At $m=9$, the orbit decomposition produces three types:
Cosmos (72 elements, period 24), Satellite (8 elements, period 8), and
Singularity (the fixed point $(9,9)$).
The canonical applied modulus is $m=24$ (the Pisano period of $m=9$),
but the orbit-type label is determined by $m=9$ arithmetic.

The central methodological principle is the \emph{observer-projection firewall}:
the EEG signals enter the QA algebraic layer exactly once (through a fixed
discretisation map), produce a discrete 4-tuple, and that 4-tuple is used
as a regression feature.  The seizure labels do not influence the discretisation.

% ====================================================================
\section{Methods}\label{sec:methods}

\subsection{Dataset}

CHB-MIT Scalp EEG Database~\cite{chbmit}: 10 patients
(chb01--chb07, chb11, chb18, chb20),
23 EEG channels, 256\,Hz sampling rate.  Seizure and non-seizure windows
are extracted as 10-second epochs; classes are balanced by undersampling
the majority class.  This patient subset was determined by data availability
at analysis time; patients chb08--chb10, chb12--chb17, and chb19, chb21--chb23
are held out and constitute a clean replication set.

\subsection{Baseline features}

Delta-band (1--4\,Hz) power per channel, averaged across the 23 channels.
This is the spectral baseline: a low-complexity feature that captures the
well-known delta-band increase during seizure onset.

\subsection{QA observer projection (Observer 3 / topographic)}

For each 10-second window:
\begin{enumerate}[label=(\arabic*)]
  \item Compute the 23-channel power spectrogram (Welch, 2\,s segments).
  \item Apply $k$-means ($k=9$) to the channel$\times$frequency matrix to
        identify spatial cluster centroids.
  \item Project each centroid to its two dominant PCA components $(c_1, c_2)$.
  \item Map to QA state: $b = (\lfloor |c_1| \rfloor \bmod 9) + 1$,
        $e = (\lfloor |c_2| \rfloor \bmod 9) + 1$.
  \item Compute the 4-tuple $(b, e, d, a)$ with
        $d = ((b+e-1)\bmod 9)+1$ and $a = ((b+2e-1)\bmod 9)+1$.
  \item Average the 4-tuple across all $k=9$ centroids for the window.
\end{enumerate}
This produces a 4-dimensional feature vector for each window
with no free parameters other than $k$ and the PCA truncation,
both fixed before seeing seizure labels.

\subsection{Regression model}

Logistic regression with L2 penalty (scikit-learn defaults).
Baseline model: delta-band features only.
Augmented model: delta-band + QA 4-tuple features.
Per-patient train/test split: 70/30, stratified.
$\Delta R^2$ is computed as McFadden pseudo-$R^2$ of augmented minus baseline model.

\subsection{Statistical testing}

Per-patient $p$-values from likelihood-ratio test comparing augmented vs.\ baseline
models.  Fisher's method combines 10 independent $p$-values:
$-2 \sum_{i=1}^{10} \ln p_i \sim \chi^2_{20}$.

% ====================================================================
\section{Results}\label{sec:results}

\begin{table}[h]
\centering
\small
\begin{tabular}{lrrl}
\toprule
\textbf{Patient} & $\Delta R^2$ & $p$ (LR test) & \textbf{Note} \\
\midrule
chb01 & $+0.200$ & $1.9 \times 10^{-5}$ & * \\
chb02 & $+0.177$ & $1.1 \times 10^{-2}$ & * \\
chb03 & $+0.228$ & $6.6 \times 10^{-7}$ & * \\
chb04 & $+0.099$ & $2.3 \times 10^{-2}$ & * \\
chb05 & $+0.277$ & $2.0 \times 10^{-8}$ & * \\
chb06 & $+0.194$ & $1.1 \times 10^{-3}$ & * \\
chb07 & $+0.192$ & $1.0 \times 10^{-3}$ & * \\
chb11 & $+0.293$ & $7.1 \times 10^{-8}$ & * \\
chb18 & $+0.048$ & $0.107$ & n.s. \\
chb20 & $+0.396$ & $1.4 \times 10^{-9}$ & * \\
\midrule
\textbf{Mean} & $\mathbf{+0.210}$ (SE $0.029$) & \multicolumn{2}{l}{Fisher $\chi^2=208.0$, $p=2.9\times10^{-33}$} \\
\bottomrule
\end{tabular}
\caption{Per-patient $\Delta R^2$ (augmented minus baseline McFadden pseudo-$R^2$)
and likelihood-ratio $p$-value ($*$ = $p < 0.05$).
Source: \texttt{eeg\_chbmit\_scale\_results.json}.}
\label{tab:results}
\end{table}

\paragraph{Aggregate.}
9/10 patients show individually significant improvement ($p < 0.05$).
The sole exception (chb18, $\Delta R^2 = +0.048$) still shows positive gain.
Fisher combined $p = 2.9 \times 10^{-33}$ is driven by the per-patient
significance rather than any single outlier: the median per-patient
$\Delta R^2 = 0.197$, close to the mean.

\paragraph{Feature ablation (HI~2.0 classifier).}
The following ablation applies to the binary HI~2.0 classifier experiment
(F1-based, separate from the Observer~3 $\Delta R^2$ result above).
The full HI~2.0 feature set comprises 10 features:
harmonic-index scalars (HI$_{2.0}$, $H_\text{ang}$, $H_\text{rad}$, $H_\text{fam}$),
Pythagorean chromogeometric elements ($C$, $F$, $G$, $\gcd$),
and raw modular coordinates ($b$, $e$).
Patients are split into two pre-defined cohorts:
\emph{stable-positive} (consistently positive HI~2.0 gain) and
\emph{hard-negative} (unreliable or negative gain).

\begin{table}[h]
\centering
\small
\begin{tabular}{lrrrr}
\toprule
& \multicolumn{2}{c}{\textbf{Stable-positive}} & \multicolumn{2}{c}{\textbf{Hard-negative}} \\
\cmidrule(lr){2-3}\cmidrule(lr){4-5}
\textbf{Feature set} & F1 & $\Delta$F1 & F1 & $\Delta$F1 \\
\midrule
Full (all 10 features)          & 0.730 & ---    & 0.644 & --- \\
No geometry ($-C,F,G,\gcd$)     & 0.693 & $-0.036$ & 0.699 & $+0.056$ \\
No coords ($-b,e$)              & 0.718 & $-0.011$ & 0.644 & $\pm 0.000$ \\
Coords only ($b,e$)             & 0.709 & $-0.021$ & 0.740 & $+0.096$ \\
Geometry only ($C,F,G,\gcd$)    & 0.715 & $-0.015$ & 0.599 & $-0.045$ \\
\bottomrule
\end{tabular}
\caption{Feature ablation for the HI~2.0 binary classifier (mean F1 across patients
in each cohort; $\Delta$F1 relative to the full feature set, positive = better than full).
This is a separate experiment from the Observer~3 $\Delta R^2$ result in Table~\ref{tab:results}.}
\label{tab:ablation}
\end{table}

The cohort split is the key finding.
For \emph{stable-positive} patients, removing the Pythagorean geometry features ($C$, $F$, $G$)
causes the largest drop ($\Delta$F1 $= -0.036$), indicating that the chromogeometric
element structure carries signal beyond the raw coordinates.
For \emph{hard-negative} patients the pattern reverses: the geometry features actively
hurt performance, and raw coordinates alone ($b$, $e$) outperform the full feature set
by $\Delta$F1 $= +0.096$.
This suggests that the two cohorts differ not only in response to QA features generally
but in \emph{which algebraic layer} carries the signal.

% ====================================================================
\section{Limitations and Replication Plan}\label{sec:limits}

\begin{enumerate}[label=(\arabic*)]
  \item \textbf{Most patients are held out.}
        The 10-patient result uses chb01--chb07, chb11, chb18, chb20.
        The remaining 13 patients (chb08--chb10, chb12--chb17, chb19, chb21--chb23)
        have not been analysed and constitute a clean hold-out set.

  \item \textbf{The observer projection involves $k$-means.}
        $k$-means is not parameter-free (requires $k=9$ and random seed).
        A replication should fix $k$ and the seed before data access
        and test sensitivity to both.

  \item \textbf{$\Delta R^2$ and $\Delta F_1$ measure different things.}
        The topographic $\Delta R^2 = +0.21$ result uses logistic regression
        pseudo-$R^2$.  A separate binary-classifier experiment (HI 2.0 vs.\ HI 1.0)
        finds mean $\Delta F_1 = +0.014$.  These are not contradictory
        but measure different aspects of seizure prediction improvement.
        The $\Delta R^2$ measures discrimination improvement;
        $\Delta F_1$ measures classification accuracy improvement.
        Both should be reported.

  \item \textbf{No held-out time period.}
        The current analysis uses all available recordings per patient.
        A prospective hold-out (e.g.\ last 20\% of each patient's recordings)
        would test temporal generalisation.
\end{enumerate}

\paragraph{Pre-registration plan.}
A replication on patients chb11--chb23 with the following pre-committed choices:
$k=9$ centroids, random seed=42, 70/30 stratified split, L2 logistic regression,
McFadden pseudo-$R^2$ as primary outcome, likelihood-ratio test per patient,
Fisher's method for aggregation.
Success criterion: mean $\Delta R^2 > 0.05$ and at least 7/13 patients significant.

% ====================================================================
\section{Discussion}\label{sec:discussion}

The $\Delta R^2 = +0.21$ gain is large by the standards of EEG feature engineering.
Comparable gains in the seizure-prediction literature typically require
patient-specific parameter tuning; the QA projection uses no patient-specific
parameters in the algebraic layer.

The most likely mechanistic explanation is that the mod-9 orbit classification
captures a discretised version of multi-channel phase relationships that are
not well represented by univariate delta-band power.
The Fibonacci shift is sensitive to the ratio $b/e$ modulo 9,
which can align with the dominant oscillation ratio across electrode clusters.
This is an explanatory hypothesis, not a claim; testing it would require
electrode-level ablation studies beyond the scope of this paper.

The cross-domain context (Section~3 of~\cite{dale2026convergence}) suggests
that the EEG gain is not an isolated artifact:
the same orbit decomposition shows structure in planetary resonances ($p < 10^{-45}$)
and climate classification ($\chi^2 = 95$).
Whether this reflects a universal property of modular Fibonacci dynamics
or domain-specific coincidences is the central open question.

\bigskip
\noindent\textbf{Acknowledgements.}
The author thanks Ben Iverson and Dale Pond for foundational QA work.
CHB-MIT data provided by the MIT Laboratory for Computational Physiology
and PhysioNet~\cite{chbmit}.

% ====================================================================
\begin{thebibliography}{10}

\bibitem{chbmit}
A.~L. Goldberger, L.~A.~N. Amaral, L.~Glass, J.~M. Hausdorff, P.~Ch. Ivanov,
R.~G. Mark, J.~E. Mietus, G.~B. Moody, C.-K. Peng, H.~E. Stanley.
\newblock PhysioBank, PhysioToolkit, and PhysioNet: components of a new research resource for complex physiologic signals.
\newblock \emph{Circulation}, 101(23):e215--e220, 2000.

\bibitem{dale2026convergence}
W.~Dale.
\newblock Modular Fibonacci orbit classification: convergent structure across graph, planetary, neural, and climate domains.
\newblock Preprint, 2026.

\bibitem{dale2026fib}
W.~Dale.
\newblock Fibonacci preference in mean-motion resonances: evidence from solar and exoplanetary systems.
\newblock Preprint, 2026.

\end{thebibliography}

\end{document}
